\chapter{Final model comparison and conclusions}

In this final section, we consolidate the results obtained across all experimental setups, integrating the generative Gaussian models from the previous laboratory with the discriminative Logistic Regression models developed in this section. The comparison evaluates both the optimal discriminatory potential (minimum DCF) and the empirical operational cost (actual DCF) for the primary application $\pi_T = 0.1$.

\section{Summary of results}

The following table summarizes the best performance achieved by each model type. For Logistic Regression models, the results corresponding to the optimal regularization parameter $\lambda$ are reported alongside their respective calibration gap.

\begin{table}[h]
\centering
\caption{Overall model comparison ($\pi_T = 0.1$)}
\label{tab:final_comparison}
\resizebox{\textwidth}{!}{
\begin{tabular}{lcccc}
\toprule
\textbf{Model Architecture} & \textbf{Optimal $\lambda$} & \textbf{minDCF} & \textbf{actDCF} & \textbf{Calibration Gap} \\
\midrule
MVG (Standard) & - & 0.2629 & 0.3051 & 0.0422 \\
Naive Bayes & - & 0.2570 & 0.3022 & 0.0452 \\
Tied MVG & - & 0.3628 & 0.4061 & 0.0432 \\
\midrule
Standard Logistic Regression & $1.00 \times 10^{-2}$ & 0.3611 & 0.4568 & 0.0956 \\
Reduced Standard LR (1/50) & $1.00 \times 10^{1}$ & 0.3783 & 1.0000 & 0.6217 \\
Prior-Weighted Logistic Regression & $3.16 \times 10^{1}$ & 0.3620 & 1.0000 & 0.6380 \\
Quadratic Logistic Regression & $3.16 \times 10^{-2}$ & \textbf{0.2436} & 0.4972 & 0.2536 \\
\bottomrule
\end{tabular}}
\end{table}

\section{Analysis and comparative discussion}

\subsection*{Which model(s) achieve(s) the best results?}

In terms of pure discriminative power, the absolute best result is achieved by the \textbf{Quadratic Logistic Regression} model, which reaches the global minimum DCF of \textbf{0.2436}. 

Following closely behind are the unconstrained generative models: \textbf{Naive Bayes} achieves a highly competitive minDCF of \textbf{0.2570}, and the \textbf{Standard MVG} follows at \textbf{0.2629}. Conversely, all models constrained to linear boundaries (Standard LR, Prior-Weighted LR, and Tied MVG) perform significantly worse, hitting a hard performance ceiling and plateauing at a minDCF around $0.361$ to $0.362$.

\subsection*{What kind of separation rules or distribution assumptions characterize this / these model(s)?}

The top-performing models (Quadratic LR, Naive Bayes, Standard MVG) all share one critical mathematical characteristic: they compute \textbf{non-linear (quadratic) decision boundaries}. 
\begin{itemize}
    \item \textbf{Quadratic Logistic Regression:} This is a discriminative model that does not make explicit assumptions about the underlying probability density of the features. By expanding the feature space to include quadratic cross-terms, the model optimizes a hyperplane in a higher-dimensional space, which maps back to a curved, quadratic decision boundary in the original feature space.
    \item \textbf{Standard MVG and Naive Bayes:} These are generative models. They explicitly assume that the data within each class follows a Normal distribution. Because they estimate class-specific covariance matrices (allowing the geometric spread of the "Fake" and "Genuine" classes to differ), the resulting Log-Likelihood Ratio forms a quadratic decision boundary.
    \item \textbf{Linear Models (Standard LR, Tied MVG):} These models are strictly limited to linear separating hyperplanes. Tied MVG forces this mechanically by assuming both classes share the exact same covariance matrix. Standard LR achieves this by optimizing a simple linear combination of the raw features without calculating interacting cross-terms.
\end{itemize}

\subsection*{How are the results related to the characteristics of the dataset features?}

By cross-referencing these final metrics with our exploratory analysis from previous laboratories, we can definitively explain the success and failure of these models based on the physical characteristics of the dataset:

\begin{enumerate}
    \item \textbf{Covariance Mismatch (Why Linear Models Fail):} The fact that Quadratic LR ($0.2436$) and Standard MVG ($0.2629$) vastly outperform Tied MVG ($0.3628$) proves that the two classes possess fundamentally different covariance structures. As observed in Laboratory 4, Genuine Feature 1 has a massively different variance compared to Fake Feature 1. A linear model cannot capture this difference in geometric spread, condemning it to a high error rate.
    \item \textbf{Non-Gaussian Skewness (Why Quadratic LR Wins):} While MVG performs well, the Quadratic Logistic Regression outperforms it. This is directly related to Features 5 and 6. In previous analyses, we proved these features are heavily skewed and do not follow a symmetric Gaussian curve. The generative MVG model suffers because its strict Normal assumption is physically violated by the data. As a discriminative model, the Quadratic LR ignores the density distributions entirely and focuses strictly on optimizing the boundary, making it robust against the skewness of Features 5 and 6 while still exploiting their variance differences.
    \item \textbf{Feature Independence (Why Naive Bayes Excels):} The Naive Bayes model ($0.2570$) slightly outperforms the full MVG model ($0.2629$). This relates back to our Laboratory 5 correlation analysis, where we proved the features are exceptionally weakly correlated (all off-diagonal coefficients hovered near zero). Because the features are already practically independent, the Naive Bayes assumption is highly accurate. By forcing the off-diagonal covariances to zero, the Naive Bayes model acts as a powerful regularizer, eliminating parameter estimation noise and slightly improving upon the full MVG.
\end{enumerate}

\section{Final considerations on calibration}

While Quadratic Logistic Regression provides the absolute best discriminatory potential (lowest minDCF), it suffers from a significant \textbf{calibration gap} of $0.2536$, yielding an actual operational DCF of $0.4972$. This massive loss of calibration is driven by the $L_2$ penalty term ($\lambda = 3.16 \times 10^{-2}$) compressing the weights and squashing the output scores. 

In stark contrast, the generative models prove to be vastly superior in native calibration. The Standard MVG and Naive Bayes models both exhibit tiny calibration gaps of just $0.0422$ and $0.0452$. Because generative models derive their scores from mathematically rigorous probability densities rather than regularized geometric distances, their outputs naturally behave as true Log-Likelihood Ratios. 

This reinforces a fundamental conclusion: while highly complex, regularized discriminative models (like Quadratic LR) can locate superior classification boundaries, their raw scores cannot be reliably thresholded out-of-the-box. Deploying such a model in a real-world security application requires a dedicated, post-processing score calibration stage. Until such calibration is applied, the generative \textbf{Naive Bayes} model remains the most reliable, "ready-to-deploy" solution for this dataset.